% Lecture 12 - Page 2 (3-Jan-1995).
% Fragment for % Lecture 12 - Page 2 (3-Jan-1995).
% Fragment for % Lecture 12 - Page 2 (3-Jan-1995).
% Fragment for % Lecture 12 - Page 2 (3-Jan-1995).
% Fragment for \input{courses/pages/curs12-pagina2-en}.
% Faithful translation, retaining manuscript sequence and notation.
% No preamble or forced page break.

\subsection*{Lecture 12 - Page 2}

\noindent\underline{Example.}\quad
$(\mathbb{Z},\cdot,1)\xrightarrow{\ f\ }(\mathbb{Z}_n,\cdot,\hat 1)$.

\noindent $f(e)=\hat e$.

\bigskip
\noindent
$(M_2(\mathbb{Z}),\cdot,I_2)\longrightarrow(\mathbb{Z},\cdot,1)$,
\qquad $A\longmapsto\det A$.
% The manuscript initially writes R in M_2(R) and in the codomain,
% with Z written underneath both symbols (an apparent correction to Z).
% Z is used here provisionally; verify against the original photograph.

\bigskip
\noindent\underline{Remark.}\quad
There also exist functions satisfying (1) but not (2)!
% The manuscript says, abbreviated: "(exists) si fct. care au (1) dar nu au (2)!".

\bigskip
\noindent $\ker f=R_f=\{(x,y)\in M\times M\mid f(x)=f(y)\}$.

\noindent $R_f$ is a congruence on $M$.

\bigskip
\noindent\underline{Proof.}\quad $R_f$ is an equivalence relation $\sim$.

\medskip
\noindent $x\,R_f\,y\ \Longrightarrow$
\[
\begin{cases}
 zx\,R_f\,zy,\\
 xz\,R_f\,yz,
\end{cases}
\qquad (\forall)\ z\in M.
\]

\noindent
$f(xz)=f(x)f(z)=f(y)f(z)=f(yz)\ \Longrightarrow\ xz\,R_f\,yz$.

\medskip
\noindent $\Longrightarrow R_f$ is compatible with ``$\cdot$''.
% The manuscript explicitly verifies right compatibility and states
% left compatibility, without a separate calculation for the latter.

\bigskip
\noindent\fbox{T.}\quad
\underline{Factorisation theorem for monoid homomorphisms}

\noindent (Fundamental isomorphism theorem for monoid homomorphisms.)

\bigskip
\noindent Let $f:M\longrightarrow M'$ be a monoid homomorphism.

\noindent $R\subseteq M\times M$ is a congruence on $M$ such that $R\subseteq R_f$.

\bigskip
\noindent\textbf{(1)}
\[
\begin{tikzcd}[column sep=large,row sep=large]
M \arrow[r,"f"] \arrow[d,"\varphi_R"'] & M' \\
M/R \arrow[ur,dashed,"f^{*}"'] &
\end{tikzcd}
\]
\noindent $\Longrightarrow\ (\exists!)\ f^{*}:M/R\longrightarrow M'$ such that
$f^{*}\circ\varphi_R=f$.
% The quotient-map label is difficult to distinguish (f_R or phi_R)
% in the original; \varphi_R is used as a provisional reading.

\bigskip
\noindent\textbf{(2)}\quad $f^{*}$ injective $\Longleftrightarrow R=R_f$.

\medskip
\noindent\textbf{(3)}\quad $f^{*}$ surjective $\Longleftrightarrow f$ surjective.

\medskip
\noindent $f^{*}$ bijective $\Longleftrightarrow$\,\footnote{In the manuscript, a curved arrow connects this equivalence to conditions (2) and (3), without writing out the right-hand side again.}

\bigskip
\noindent\fbox{T.}\quad (Fundamental isomorphism theorem).\quad
$f:M\longrightarrow M'$ a monoid homomorphism
\[
\Longrightarrow\quad \operatorname{Im} f\cong\frac{M}{R_f}.
\]
\noindent\hfill $f(M)\leq M'$.
% "f(M) <= M'" appears beneath "Im f" in the manuscript.
% The page ends here.
.
% Faithful translation, retaining manuscript sequence and notation.
% No preamble or forced page break.

\subsection*{Lecture 12 - Page 2}

\noindent\underline{Example.}\quad
$(\mathbb{Z},\cdot,1)\xrightarrow{\ f\ }(\mathbb{Z}_n,\cdot,\hat 1)$.

\noindent $f(e)=\hat e$.

\bigskip
\noindent
$(M_2(\mathbb{Z}),\cdot,I_2)\longrightarrow(\mathbb{Z},\cdot,1)$,
\qquad $A\longmapsto\det A$.
% The manuscript initially writes R in M_2(R) and in the codomain,
% with Z written underneath both symbols (an apparent correction to Z).
% Z is used here provisionally; verify against the original photograph.

\bigskip
\noindent\underline{Remark.}\quad
There also exist functions satisfying (1) but not (2)!
% The manuscript says, abbreviated: "(exists) si fct. care au (1) dar nu au (2)!".

\bigskip
\noindent $\ker f=R_f=\{(x,y)\in M\times M\mid f(x)=f(y)\}$.

\noindent $R_f$ is a congruence on $M$.

\bigskip
\noindent\underline{Proof.}\quad $R_f$ is an equivalence relation $\sim$.

\medskip
\noindent $x\,R_f\,y\ \Longrightarrow$
\[
\begin{cases}
 zx\,R_f\,zy,\\
 xz\,R_f\,yz,
\end{cases}
\qquad (\forall)\ z\in M.
\]

\noindent
$f(xz)=f(x)f(z)=f(y)f(z)=f(yz)\ \Longrightarrow\ xz\,R_f\,yz$.

\medskip
\noindent $\Longrightarrow R_f$ is compatible with ``$\cdot$''.
% The manuscript explicitly verifies right compatibility and states
% left compatibility, without a separate calculation for the latter.

\bigskip
\noindent\fbox{T.}\quad
\underline{Factorisation theorem for monoid homomorphisms}

\noindent (Fundamental isomorphism theorem for monoid homomorphisms.)

\bigskip
\noindent Let $f:M\longrightarrow M'$ be a monoid homomorphism.

\noindent $R\subseteq M\times M$ is a congruence on $M$ such that $R\subseteq R_f$.

\bigskip
\noindent\textbf{(1)}
\[
\begin{tikzcd}[column sep=large,row sep=large]
M \arrow[r,"f"] \arrow[d,"\varphi_R"'] & M' \\
M/R \arrow[ur,dashed,"f^{*}"'] &
\end{tikzcd}
\]
\noindent $\Longrightarrow\ (\exists!)\ f^{*}:M/R\longrightarrow M'$ such that
$f^{*}\circ\varphi_R=f$.
% The quotient-map label is difficult to distinguish (f_R or phi_R)
% in the original; \varphi_R is used as a provisional reading.

\bigskip
\noindent\textbf{(2)}\quad $f^{*}$ injective $\Longleftrightarrow R=R_f$.

\medskip
\noindent\textbf{(3)}\quad $f^{*}$ surjective $\Longleftrightarrow f$ surjective.

\medskip
\noindent $f^{*}$ bijective $\Longleftrightarrow$\,\footnote{In the manuscript, a curved arrow connects this equivalence to conditions (2) and (3), without writing out the right-hand side again.}

\bigskip
\noindent\fbox{T.}\quad (Fundamental isomorphism theorem).\quad
$f:M\longrightarrow M'$ a monoid homomorphism
\[
\Longrightarrow\quad \operatorname{Im} f\cong\frac{M}{R_f}.
\]
\noindent\hfill $f(M)\leq M'$.
% "f(M) <= M'" appears beneath "Im f" in the manuscript.
% The page ends here.
.
% Faithful translation, retaining manuscript sequence and notation.
% No preamble or forced page break.

\subsection*{Lecture 12 - Page 2}

\noindent\underline{Example.}\quad
$(\mathbb{Z},\cdot,1)\xrightarrow{\ f\ }(\mathbb{Z}_n,\cdot,\hat 1)$.

\noindent $f(e)=\hat e$.

\bigskip
\noindent
$(M_2(\mathbb{Z}),\cdot,I_2)\longrightarrow(\mathbb{Z},\cdot,1)$,
\qquad $A\longmapsto\det A$.
% The manuscript initially writes R in M_2(R) and in the codomain,
% with Z written underneath both symbols (an apparent correction to Z).
% Z is used here provisionally; verify against the original photograph.

\bigskip
\noindent\underline{Remark.}\quad
There also exist functions satisfying (1) but not (2)!
% The manuscript says, abbreviated: "(exists) si fct. care au (1) dar nu au (2)!".

\bigskip
\noindent $\ker f=R_f=\{(x,y)\in M\times M\mid f(x)=f(y)\}$.

\noindent $R_f$ is a congruence on $M$.

\bigskip
\noindent\underline{Proof.}\quad $R_f$ is an equivalence relation $\sim$.

\medskip
\noindent $x\,R_f\,y\ \Longrightarrow$
\[
\begin{cases}
 zx\,R_f\,zy,\\
 xz\,R_f\,yz,
\end{cases}
\qquad (\forall)\ z\in M.
\]

\noindent
$f(xz)=f(x)f(z)=f(y)f(z)=f(yz)\ \Longrightarrow\ xz\,R_f\,yz$.

\medskip
\noindent $\Longrightarrow R_f$ is compatible with ``$\cdot$''.
% The manuscript explicitly verifies right compatibility and states
% left compatibility, without a separate calculation for the latter.

\bigskip
\noindent\fbox{T.}\quad
\underline{Factorisation theorem for monoid homomorphisms}

\noindent (Fundamental isomorphism theorem for monoid homomorphisms.)

\bigskip
\noindent Let $f:M\longrightarrow M'$ be a monoid homomorphism.

\noindent $R\subseteq M\times M$ is a congruence on $M$ such that $R\subseteq R_f$.

\bigskip
\noindent\textbf{(1)}
\[
\begin{tikzcd}[column sep=large,row sep=large]
M \arrow[r,"f"] \arrow[d,"\varphi_R"'] & M' \\
M/R \arrow[ur,dashed,"f^{*}"'] &
\end{tikzcd}
\]
\noindent $\Longrightarrow\ (\exists!)\ f^{*}:M/R\longrightarrow M'$ such that
$f^{*}\circ\varphi_R=f$.
% The quotient-map label is difficult to distinguish (f_R or phi_R)
% in the original; \varphi_R is used as a provisional reading.

\bigskip
\noindent\textbf{(2)}\quad $f^{*}$ injective $\Longleftrightarrow R=R_f$.

\medskip
\noindent\textbf{(3)}\quad $f^{*}$ surjective $\Longleftrightarrow f$ surjective.

\medskip
\noindent $f^{*}$ bijective $\Longleftrightarrow$\,\footnote{In the manuscript, a curved arrow connects this equivalence to conditions (2) and (3), without writing out the right-hand side again.}

\bigskip
\noindent\fbox{T.}\quad (Fundamental isomorphism theorem).\quad
$f:M\longrightarrow M'$ a monoid homomorphism
\[
\Longrightarrow\quad \operatorname{Im} f\cong\frac{M}{R_f}.
\]
\noindent\hfill $f(M)\leq M'$.
% "f(M) <= M'" appears beneath "Im f" in the manuscript.
% The page ends here.
.
% Faithful translation, retaining manuscript sequence and notation.
% No preamble or forced page break.

\subsection*{Lecture 12 - Page 2}

\noindent\underline{Example.}\quad
$(\mathbb{Z},\cdot,1)\xrightarrow{\ f\ }(\mathbb{Z}_n,\cdot,\hat 1)$.

\noindent $f(e)=\hat e$.

\bigskip
\noindent
$(M_2(\mathbb{Z}),\cdot,I_2)\longrightarrow(\mathbb{Z},\cdot,1)$,
\qquad $A\longmapsto\det A$.
% The manuscript initially writes R in M_2(R) and in the codomain,
% with Z written underneath both symbols (an apparent correction to Z).
% Z is used here provisionally; verify against the original photograph.

\bigskip
\noindent\underline{Remark.}\quad
There also exist functions satisfying (1) but not (2)!
% The manuscript says, abbreviated: "(exists) si fct. care au (1) dar nu au (2)!".

\bigskip
\noindent $\ker f=R_f=\{(x,y)\in M\times M\mid f(x)=f(y)\}$.

\noindent $R_f$ is a congruence on $M$.

\bigskip
\noindent\underline{Proof.}\quad $R_f$ is an equivalence relation $\sim$.

\medskip
\noindent $x\,R_f\,y\ \Longrightarrow$
\[
\begin{cases}
 zx\,R_f\,zy,\\
 xz\,R_f\,yz,
\end{cases}
\qquad (\forall)\ z\in M.
\]

\noindent
$f(xz)=f(x)f(z)=f(y)f(z)=f(yz)\ \Longrightarrow\ xz\,R_f\,yz$.

\medskip
\noindent $\Longrightarrow R_f$ is compatible with ``$\cdot$''.
% The manuscript explicitly verifies right compatibility and states
% left compatibility, without a separate calculation for the latter.

\bigskip
\noindent\fbox{T.}\quad
\underline{Factorisation theorem for monoid homomorphisms}

\noindent (Fundamental isomorphism theorem for monoid homomorphisms.)

\bigskip
\noindent Let $f:M\longrightarrow M'$ be a monoid homomorphism.

\noindent $R\subseteq M\times M$ is a congruence on $M$ such that $R\subseteq R_f$.

\bigskip
\noindent\textbf{(1)}
\[
\begin{tikzcd}[column sep=large,row sep=large]
M \arrow[r,"f"] \arrow[d,"\varphi_R"'] & M' \\
M/R \arrow[ur,dashed,"f^{*}"'] &
\end{tikzcd}
\]
\noindent $\Longrightarrow\ (\exists!)\ f^{*}:M/R\longrightarrow M'$ such that
$f^{*}\circ\varphi_R=f$.
% The quotient-map label is difficult to distinguish (f_R or phi_R)
% in the original; \varphi_R is used as a provisional reading.

\bigskip
\noindent\textbf{(2)}\quad $f^{*}$ injective $\Longleftrightarrow R=R_f$.

\medskip
\noindent\textbf{(3)}\quad $f^{*}$ surjective $\Longleftrightarrow f$ surjective.

\medskip
\noindent $f^{*}$ bijective $\Longleftrightarrow$\,\footnote{In the manuscript, a curved arrow connects this equivalence to conditions (2) and (3), without writing out the right-hand side again.}

\bigskip
\noindent\fbox{T.}\quad (Fundamental isomorphism theorem).\quad
$f:M\longrightarrow M'$ a monoid homomorphism
\[
\Longrightarrow\quad \operatorname{Im} f\cong\frac{M}{R_f}.
\]
\noindent\hfill $f(M)\leq M'$.
% "f(M) <= M'" appears beneath "Im f" in the manuscript.
% The page ends here.
